\documentclass[14 pt]{extarticle}

	\usepackage[frenchb]{babel}
	\usepackage[utf8]{inputenc}  
	\usepackage[T1]{fontenc}
	\usepackage{amssymb}
	\usepackage[mathscr]{euscript}
	\usepackage{stmaryrd}
	\usepackage{amsmath}
	\usepackage{tikz}
	\usepackage[all,cmtip]{xy}
	\usepackage{amsthm}
	\usepackage{varioref}
	\usepackage{geometry}
	\geometry{a4paper}
	\usepackage{lmodern}
	\usepackage{hyperref}
	\usepackage{array}
	 \usepackage{fancyhdr}
	 \usepackage{float}
\renewcommand{\theenumi}{\alph{enumi})}
	\pagestyle{fancy}
	\theoremstyle{plain}
	\fancyfoot[C]{} 
	\fancyhead[L]{Contrôle}
	\fancyhead[R]{3 février 2026}\geometry{
 a4paper,
 total={190mm,257mm},
 left=10mm,
 top=20mm,
 }
	
	
	\title{Contrôle Chapitre 2}
	\date{}
	\begin{document}

\begin{center}{\Large Interrogation Chapitre 6}\\ 
 \end{center}
 Nom : \\
 Prénom : \\
 \subsection*{Exercice 1 (3 points)}
 
 Rappelez la définition de la médiatrice d'un segment $[AB]$. 
  
\subsection*{Exercice 2 (10 points)}

\begin{enumerate}
\item Tracer les médiatrices des segments $[FS]$, $[IS]$ et $[IA]$. 
\item Tracer le cercle passant par $F$, $I$ et $S$. 
\item Existe-t-il un cercle passant par les quatre points ? Justifier votre réponse \textbf{sur votre copie.}
\end{enumerate}
\includegraphics[scale=1]{Exercice1}


\subsection*{Exercice 3 (2 points)}
Énoncez le théorème de la médiatrice. 

\subsection*{Exercice 4 (5 points)}
\includegraphics[scale=1]{Exercice2}
\begin{enumerate}
\item Que représente la droite $(d)$ pour le segment $[JL]$ ? Justifier soigneusement votre réponse. 
\item Démontrer soigneusement que le triangle $KJL$ est isocèle en $K$.
\end{enumerate}

 	\end{document}
